\section{Theory and Methodology}\label{sec:theory}

\subsection{Setup and assumptions}\label{sec:assumptions}

We consider multiclass classification with inputs $x\in\mathcal{X}$ and labels
$y\in\{1,\dots,K\}$. A training sample $S=\{(x_i,y_i)\}_{i=1}^{n}$ is drawn
\emph{i.i.d.} from an unknown distribution $\mathcal{D}$. A stochastic predictor
is a distribution $Q$ over weight vectors $w$; drawing $w\sim Q$ and predicting
$f_w(x)$ defines the \emph{Gibbs classifier}. For a loss $\ell$ taking values in
$[0,1]$ the Gibbs risk and its empirical counterpart on a set $T$ are
\begin{equation}\label{eq:risk}
R(Q)=\mathbb{E}_{w\sim Q}\,\mathbb{E}_{(x,y)\sim\mathcal{D}}\,\ell(f_w(x),y),
\qquad
\hat{R}_T(Q)=\mathbb{E}_{w\sim Q}\,\frac{1}{|T|}\sum_{(x,y)\in T}\ell(f_w(x),y).
\end{equation}
A certificate bounds the population risk using the empirical quantity.
PAC-Bayes bounds require (i) i.i.d.\ data, (ii) a loss bounded in $[0,1]$, and
(iii) a prior $P$ chosen \emph{independently of the sample used to evaluate the
empirical risk in the bound}. The $0$--$1$ loss is bounded directly; the
cross-entropy is not, so wherever it enters a bound we use the standard bounded
surrogate of \textcite{dziugaite2017}, clamping log-probabilities at
$\log p_{\min}$ and rescaling by $1/\log(1/p_{\min})$, with $p_{\min}=10^{-5}$.

\subsection{The PAC-Bayes-kl certificate}\label{sec:pbkl}

All formal certificates in this dissertation are instances of the
PAC-Bayes-kl inequality \parencite{maurer2004,seeger2002}, the tightest of the
classical PAC-Bayes statements. Writing
$\mathrm{kl}(q\|p)=q\log\frac{q}{p}+(1-q)\log\frac{1-q}{1-p}$ for the binary KL,
it states that for any prior $P$ independent of the bound set $S_{\mathrm b}$ of
size $n_{\mathrm b}$, simultaneously for all posteriors $Q$, with probability at
least $1-\delta$ over the draw of $S_{\mathrm b}$,
\begin{equation}\label{eq:pbkl}
\mathrm{kl}\!\big(\hat{R}_{S_{\mathrm b}}(Q)\,\big\|\,R(Q)\big)\;\le\;
\frac{\mathrm{KL}(Q\|P)+\log\frac{2\sqrt{n_{\mathrm b}}}{\delta}}{n_{\mathrm b}}.
\end{equation}
The certificate is obtained by inverting~\eqref{eq:pbkl}: with
$\mathrm{kl}^{-1}(u;c)=\sup\{v\in[u,1):\mathrm{kl}(u\|v)\le c\}$,
\begin{equation}\label{eq:invert}
R(Q)\;\le\;\mathrm{kl}^{-1}\!\Big(\hat{R}_{S_{\mathrm b}}(Q)\,;\,
\tfrac{1}{n_{\mathrm b}}\big(\mathrm{KL}(Q\|P)+\log\tfrac{2\sqrt{n_{\mathrm b}}}{\delta}\big)\Big),
\end{equation}
which has no closed form and is evaluated by bisection. Two quantities govern it:
the empirical Gibbs risk $\hat{R}_{S_{\mathrm b}}(Q)$ and the complexity term
$\mathrm{KL}(Q\|P)/n_{\mathrm b}$. The sample size in~\eqref{eq:pbkl} is
$n_{\mathrm b}$, the bound set on which the empirical risk is measured, not the
full dataset; Section~\ref{sec:partition} returns to why that matters.

\subsection{What the certificate covers: the Gibbs classifier}\label{sec:scope}

The object bounded by~\eqref{eq:invert} is the Gibbs classifier, which freshly
samples $w\sim Q$ for each prediction --- rarely what a practitioner deploys,
since sampling per example is wasteful and makes the prediction random. We
therefore report three predictors: the \emph{stochastic} (Gibbs) one, which the
certificate controls; the \emph{posterior-mean} predictor $\mathbb{E}_Q[w]$, the
natural deterministic deployment; and the \emph{finite ensemble} averaging the
softmax outputs of $L$ draws. Reporting all three shows how far the certified
quantity sits from what would be deployed.

The posterior-mean and ensemble errors are functions of $Q$ but are \emph{not}
controlled by~\eqref{eq:pbkl}: a low ordering on a finite test set does
not transfer the Gibbs certificate to a deterministic predictor. A formal
guarantee for the majority vote is obtainable in principle from the Gibbs risk
and its variance via the C-bound \parencite{lacasse2006,germain2015}, but it is a
different, generally looser statement that we do not compute here.

\subsection{From finite Monte-Carlo to a certificate: two failure
probabilities}\label{sec:confidence}

The empirical Gibbs risk $\hat{R}_{S_{\mathrm b}}(Q)$ in~\eqref{eq:invert} is an
expectation over $w\sim Q$ and cannot be computed exactly; it is estimated from
$m$ independent Monte-Carlo replicates. This introduces a \emph{second} source of
randomness with its own failure probability, which must be accounted for
separately from the PAC-Bayes one, so we split the confidence budget in two and
combine the parts by a union bound. Let $\bar{R}$ be the mean $0$--$1$
error over the $m$ replicates and the $n_{\mathrm b}$ bound examples. A replicate
scores one draw $w\sim Q$ on the whole bound set, or, where memory forces that set
into several batches, one independent draw per batch. Either way a replicate is an
unbiased $[0,1]$-valued estimate of $\hat{R}_{S_{\mathrm b}}(Q)$ and the $m$ of
them are i.i.d., which is what~\eqref{eq:mc} needs: the count entering the slack is
$m$, not the number of forward passes. First, inverting
the binary KL with the Monte-Carlo slack controls the sampling error with
probability $\ge 1-\delta_{\mathrm{MC}}$ \parencite{langford2001}:
\begin{equation}\label{eq:mc}
\hat{R}_{S_{\mathrm b}}(Q)\;\le\;\overline{R}^{+}:=\mathrm{kl}^{-1}\!\Big(\bar{R};\,
\tfrac{1}{m}\log\tfrac{2}{\delta_{\mathrm{MC}}}\Big).
\end{equation}
Second, substituting $\overline{R}^{+}$ for $\hat{R}_{S_{\mathrm b}}(Q)$
in~\eqref{eq:invert} at confidence $\delta_{\mathrm{PAC}}$ yields the
population certificate. By a union bound over the two events that can fail,
the final certificate holds with probability at least
\begin{equation}\label{eq:joint}
1-\delta_{\mathrm{MC}}-\delta_{\mathrm{PAC}}.
\end{equation}
It is the sum in~\eqref{eq:joint}, not either $\delta$ alone, that we report.
With $\delta_{\mathrm{MC}}=\delta_{\mathrm{PAC}}=0.005$ every reported
certificate holds with probability at least $0.99$; reusing a single $\delta=0.01$
across both inversions, as the reference does, would double-count one failure
probability and overstate the confidence. A third parameter
$\delta_{\mathrm{train}}=0.025$ enters only the training objective
(Section~\ref{sec:objectives}) and never the reported certificate; all three are
recorded separately for every run. The headline certificate throughout is the
$0$--$1$ one; the cross-entropy certificate is auxiliary, and we do not claim the
two hold \emph{simultaneously} at $0.99$, which would need a further union split.
Figure~\ref{fig:chain} traces the whole computation.

\begin{figure}[t]\centering
\includegraphics[width=\linewidth]{fig_certificate_chain.pdf}
\caption{The certificate computation, end to end. Two distinct corrections
stand between the raw Monte-Carlo estimate and the reported number, and the two
leader lines mark them as the only two events that can fail: their union is what
gives the $0.99$. The dashed branch from the prior records that $\sigma_0$ was
chosen by inspecting $S_{\mathrm b}$, which is paid for by splitting
$\delta_{\mathrm{PAC}}$ across the candidate grid rather than by a separate
failure probability. Bottom right: all three predictors come from the same $Q$,
but only the Gibbs one is certified, and the crossed lines mark that the
guarantee does not transfer to the other two.}\label{fig:chain}
\end{figure}

\paragraph{Choosing the prior spends part of the same budget.} The prior
scale $\sigma_0$ was not fixed in advance: it was settled after comparing
certificates computed on $S_{\mathrm b}$ (Section~\ref{sec:protocol}).
Inequality~\eqref{eq:pbkl} holds for a prior fixed independently of
$S_{\mathrm b}$, simultaneously over all posteriors; it does not licence
inspecting $S_{\mathrm b}$ and then picking one prior out of several for free.
The repair is one more union bound: applying~\eqref{eq:pbkl} to each of $K$
candidate priors at $\delta_{\mathrm{PAC}}/K$ makes the guarantee simultaneous
over all of them, so whichever is adopted is covered, at a cost of $\log K$ in
the numerator of the slack.

The obvious candidate set is the wrong one. Counting the scales we compared, the
three probed and the one adopted, gives $K=4$; but $0.03$ was picked \emph{after}
those curves were seen, so that set is itself a function of $S_{\mathrm b}$ and a
union over it covers nothing. We union instead over
$\Sigma=\{0.001,0.002,\dots,0.999\}$, every three-decimal scale in the unit
interval, so $K=999$. The price, $\log 999$ nats against $\log 4$, raises the
headline certificates by about $0.0005$ and the smallest bound set in the study,
the $10\%$-data cell at $n_{\mathrm b}=3{,}000$, by $0.004$; no comparison here
turns on that. We apply it to every cell, conservatively so for Fashion-MNIST and
CIFAR-10, whose bound sets were never consulted in choosing $\sigma_0$.

The union is
valid for any $\Sigma$ fixed independently of $S_{\mathrm b}$, and then covers the
adopted scale however it was picked from $\Sigma$. What we cannot show is that our
selection was confined to $\Sigma$ in advance: $\Sigma$ was written down
afterwards, and a procedure returning $0.0325$ would not have been covered. The
certificates therefore hold under the assumption that only a three-decimal scale
could have been produced, which we believe but cannot demonstrate from the
artefact. Pre-registration is one way to close that, but not the only one: it
would equally suffice to show the candidate set fixed in advance by the
configuration space or by the search procedure itself. Ours was not, so we carry
it as a threat to validity (Section~\ref{sec:threats}) rather than as a solved
problem.

The posterior needs no such treatment, since~\eqref{eq:pbkl} is already
simultaneous over all $Q$: the learning rate, the epoch count and everything else
shaping $Q$ may depend on the data freely. That licence covers the PAC-Bayes
event only, not the finite-Monte-Carlo step~\eqref{eq:mc}, which holds for the
$Q$ evaluated rather than simultaneously over candidates. Choosing a checkpoint
by comparing finite-$m$ certificates would therefore need its own correction; we
do not, computing each certificate once on the final-epoch posterior.

\subsection{Data partition: \texorpdfstring{$S$, $S_0$ and $S\setminus
S_0$}{S, S0 and S minus S0}}\label{sec:partition}

A learnt (data-dependent) prior is built by training on part of the data, which
is legitimate only if the bound set is kept independent of the prior. From the
full training set we
select a subset $S$, which is the whole training set unless we are studying data
scarcity. For a learnt prior we split $S$ into a \emph{prior subset}
$S_0\subset S$ (half of $S$) used to train the prior, and the disjoint \emph{bound
subset} $S_{\mathrm b}:=S\setminus S_0$ on which the empirical risk entering the
certificate is measured. The posterior $Q$ may be trained on all of $S$; the prior
$P$ depends only on $S_0$ and is therefore independent of $S_{\mathrm b}$,
as~\eqref{eq:pbkl} requires. Consequently
\begin{equation}\label{eq:counts}
n_{\mathrm{post}}=|S|,\qquad n_{\mathrm b}=|S_{\mathrm b}|.
\end{equation}
The second equality in~\eqref{eq:counts} is the one the audit of
Section~\ref{sec:pipeline} restores: it is the bound subset, not the dataset, that
supplies the denominator.
For a data-free (random) prior there is no $S_0$, the prior is independent of all
of $S$, and $n_{\mathrm b}=|S|$. On the full MNIST training set with an even
split this gives $n_{\mathrm{post}}=60{,}000$ and $n_{\mathrm b}=30{,}000$
($50{,}000$ and $25{,}000$ for CIFAR-10); putting $60{,}000$ in the denominator
instead, as a na\"ive reading of the partition suggests, understates the
complexity term and returns a certificate that is too tight. We therefore derive
$n_{\mathrm b}$ from explicit bound-subset membership and verify
$S_0\cap S_{\mathrm b}=\varnothing$ and $S_0\cup S_{\mathrm b}=S$ before
every run.

\begin{figure}[t]\centering
\includegraphics[width=\linewidth]{fig_partition.pdf}
\caption{Data partition and the two prior cases. A uniformly random subset $S$ is
drawn from the full training set, and in Case~1 split again into a prior subset
$S_0$ and a bound subset $S_{\mathrm b}=S\setminus S_0$. Both draws are made by
permutation alone and never read a label, which is what leaves $P$ admissible;
Figure~\ref{fig:flow} gives the dependence this buys. The certificate's
denominator is the size of whichever set the empirical risk is measured on, so it
is $n_{\mathrm b}=|S_{\mathrm b}|$ with a learnt prior and $n_{\mathrm b}=|S|$
with a data-free one --- the distinction the sample-count correction of
Section~\ref{sec:pipeline} restores.}\label{fig:partition}
\end{figure}

Figure~\ref{fig:partition} fixes which examples go where; what the bound
constrains is which subset may \emph{inform} which quantity
(Figure~\ref{fig:flow}).

\begin{figure}[t]\centering
\includegraphics[width=0.92\linewidth]{fig_information_flow.pdf}
\caption{Permitted and forbidden information flow. The posterior may be informed
by every part of $S$; the prior may be informed only by $S_0$. The one forbidden
edge is $S_{\mathrm b}\to P$, and the note records why disjoint index sets are
not by themselves enough to rule it out: a class-stratified split sets $S_0$'s
quotas from label counts taken over all of $S$, so the prior can depend on the
bound set's labels without sharing a single example with it. That is the failure
the label-independent partition of Section~\ref{sec:pipeline}
removes.}\label{fig:flow}
\end{figure}

\subsection{Training objectives}\label{sec:objectives}

The posterior is trained by PAC-Bayes-with-Backprop \parencite{perezortiz2021}:
the weight distribution $Q$ (a diagonal Gaussian, mean and per-weight scale) is
optimised by gradient descent on a differentiable objective derived from a
PAC-Bayes bound, through the reparameterisation trick \parencite{blundell2015}.
Whether that objective is itself a bound depends on $\beta$, settled below. We compare three
bound-derived objectives and one accuracy-oriented baseline, all on
$n=n_{\mathrm{post}}$ training points with bounded empirical cross-entropy
$\hat{R}(Q)$ and confidence $\delta_{\mathrm{train}}$:
Writing $\kappa=\beta\,\mathrm{KL}(Q\|P)+\log\big(2\sqrt{n}/\delta_{\mathrm{train}}\big)$
for the complexity numerator shared by the first three:
\begin{itemize}
\item \textbf{$f_{\text{classic}}$} (McAllester/Pinsker):
  $f=\hat{R}(Q)+\sqrt{\kappa/(2n)}$.
\item \textbf{$f_{\lambda}$} \parencite{thiemann2017}: for a learnable
  $\lambda\in(0,2)$, optimised by alternating minimisation over $Q$ and $\lambda$,
  \[
  f=\frac{\hat{R}(Q)}{1-\lambda/2}+\frac{\kappa}{n\,\lambda\,(1-\lambda/2)}.
  \]
\item \textbf{$f_{\text{quad}}$}: the quadratic relaxation of PAC-Bayes-kl
  \parencite{perezortiz2021}, which with $\rho=\kappa/(2n)$ is
  \[
  f=\Big(\sqrt{\hat{R}(Q)+\rho}+\sqrt{\rho}\Big)^{2}.
  \]
  It is reported there as the tightest of the three with a data-dependent prior,
  though in our runs $f_{\text{classic}}$ is tighter (Section~\ref{sec:main}).
\item \textbf{$f_{\text{bbb}}$} (baseline): the Bayes-by-Backprop objective
  \parencite{blundell2015}, $f=\hat{R}(Q)+\beta\,\mathrm{KL}(Q\|P)/n$. It has the
  ELBO's form but is not the ELBO: $\hat{R}(Q)$ here is the clamped, rescaled
  cross-entropy of Section~\ref{sec:assumptions} rather than a log-likelihood, and
  $\beta<1$ down-weights the KL. It trades that KL against the empirical risk with
  a calibration that does not target the reported certificate, which is what makes
  it an accuracy-oriented control.
\end{itemize}
The KL penalty $\beta$ is carried by all four objectives, not only the baseline,
and following the reference's configuration we set $\beta=0.1$ throughout. This
makes each of the first three a down-weighted surrogate rather than the literal
bound it is named after, so the value each optimises is not itself a valid
certificate. That does not affect what we report, because the objective only
shapes $Q$: the reported certificate is always the \emph{same} PAC-Bayes-kl
inversion~\eqref{eq:invert} applied to the resulting $Q$, evaluated with the
unpenalised $\mathrm{KL}(Q\|P)$. No ordering of objectives by certificate value
is therefore fixed independently of the $Q$ they produce.

\subsection{Priors}\label{sec:priors}

We study two prior families. A \emph{data-free} (random) prior is a fixed
isotropic Gaussian at the initialisation, independent of all data: the most
distribution-independent choice, but one that typically leaves a large
$\mathrm{KL}(Q\|P)$ and a loose certificate. A \emph{data-dependent} (learnt)
prior centres $P$ at the weights of a network trained by empirical-risk
minimisation on $S_0$, with dropout to limit prior over-fitting, which moves $P$
close to well-performing posteriors and can sharply reduce the KL. The posterior
is initialised at the prior in both cases, and the legitimacy of the learnt one
rests on the independence of $P$ from $S\setminus S_0$
(Section~\ref{sec:partition}).
